% !TeX root = ../papers/nc-p3.tex
\section{二次曲面との交点による枠のKoszulフィルトレーション}
\label{sec:half-period-semisimple}
\begin{theorem}\label{thm:half-period-semisimple}
前節の枠で $e_i$ を $u_i$ の双対基底とし、$G=(q(e_i,e_j))$ とする。
$G,G^{-1}$ の対角項が全て非零とする。
このとき標数零の代数閉体で、全ての非零重み $\sum\omega_i=1$ に対し
窓は標準Koszul左右AS正則で $q$ は非零因子である。
これは九次元成分の稠密開部分の全閉点を含む有限な十分条件である。
\end{theorem}
\begin{proof}[八次元半単純環での二次双対]
$D=k^4\oplus M_2(k)$ に積を入れ、$\ell(t,Z)=\sum_i t_i-\operatorname{tr}(Z)/2$、
$H=\ker\ell$ とする。$q(x,x)=\det X(x)$ となる線型同型 $X:E\to M_2(k)$ を選ぶ。
偶・奇の一次空間を $D$ に次の四空間 $V,W$ として入れる。
\[
 x\mapsto ((u_i(x)),X(x)),\qquad y\mapsto ((v_i(y)),\operatorname{adj}X(y)).
\]
直線族の順序の交換が必要なら転置を施す。$\operatorname{adj}$ は二次行列では線型である。
両空間ともスカラー四成分への射影と行列成分への射影は同型である。
極化恒等式で $\ell(VW)=\ell(WV)=0$ であり、非退化な跡対合 $\ell(ab)$ に関して $W=V^\perp$ である。
二周期でこれらを交互に掛けて生成する次数付き代数を $C$ とする（次数ごとに $D$ の別のコピーを使う）。
二次積のスカラー成分は四つの $p_i$、行列成分の跡零部分は $L_+$ または $L_-$ の三形式である。
実際、$X\operatorname{adj}Y$ の跡零部分は交代的で、階数一の二行列に共通の核があれば零である。
従ってその三形式は一方の直線族を消す三空間を張る。逆順積では共通の像を使い他方の族となる。
$\sum c_i p_i$ の対角を $e_j$ で評価すると $c_jG_{jj}$ となるので $K$ は交代空間に交わらない。
よって二次積の像は各位置でちょうど $H$、核は中心商 $B$ の関係空間の直交である。

$V,W$ はともに $D$ の単元を稠密に含む。$H$ の右安定化単元は
スカラー単元だけである。実際 $H$ の直交直線は $1$ で張られ、跡対合は非退化だからである。
従って $H$ に比例しない二単元を掛けた二超平面は異なり、三次積の像は $D$ になる。
以後も単元を掛ければ $D$ である。従って各始点で $C_n$ は次数順に $k,V$ または $W,H,D,D,\ldots$ となる。
\end{proof}

\begin{proof}[線型イデアルの構成]
始点を固定してその一次空間を $V_0$、次を $V_1$ とする（$V,W$ のどちらから始めてもよい）。
$D$ の右イデアルを
\[
 D_{S,L}=\{(t,Z):t_i=0\ (i\notin S),\ \operatorname{im}Z\subset L\}
 \quad(S\subset\{1,2,3,4\},\ L\subset k^2)
\]
と書き、これと正次数の $C$ との次数別交叉を $J_{S,L}$ とする。次数零は零とする。
$V_p$ のスカラーが $S$ に台を持つ部分は $|S|$ 次元である。
四添字を任意に二対に分け、そのいずれかを $S$ とする。対応する行列二空間の射影直線は行列式二次曲面に含まれない。
偶・奇の行列式は対角基底で $G$ または $G^{-1}$ の形式（共通倍率を除く）で、
各端の座標行列は可逆である。直線は少なくとも一つ階数1の根を持ち、その台はちょうど $S$ である。二重根も許す。以下では行列成分を大文字で、対応する一次元を小文字で表す。

根 $x$ の行列を $X$ とする。右零化空間は $D_{S^c,\ker X}$ で、
次位置の相補割線には像が $\ker X$ の根 $y$ がちょうど一つある。
実際、相補二空間は行列跡で直交する。第一の根 $X$ と同じ二空間内の可逆な $X'$ を基底とする。
像を $\ker X$ に制限した行列二空間では $\operatorname{tr}(XY)=0$ は自動で、
$\operatorname{tr}(X'Y)=0$ は非零の一方程式なので解は直線である。
従って $x V_1$ は次元3で $H\cap D_{S,\operatorname{im}X}$ に等しく、
$xH=xD=D_{S,\operatorname{im}X}$ は次元4である。後者は $\ker(x\cdot)$ に
$S^c$ 上の単位座標があり $\ell$ が非零だからである。一次でも同じ二空間に可逆な行列があるので
$J_{S,\operatorname{im}X}=xC$、その右零化イデアルは $yC$ となる。この議論は全根と両始点に適用する。

次にこの根 $x$ と同じ二空間の可逆行列を持つ $x'$ を加える。$xC$ に対する $x'$ の右コロン
$\{b:x'b\in xC\}$ は次位置の $J_{\mathrm{all},L'}$ である。
ただし $L'=X'^{-1}\operatorname{im}X$ である。正次数の積では $xC$ は上記の交叉を全て取り、$x'$ のスカラー台は $S$ 内であるからである。
このコロンの一次元数は2、二次は5、以後は6である。
従って $(x,x')C=J_{S,k^2}$ は次数1,2,3以後で次元 $2,5,6$ となる。
実際 $(x,x')C/xC$ はコロンによる商の一次ずらしであり、各次数で交叉の次元に達する。

ここで現れた $L=L'$ の $J_{\mathrm{all},L}$ も一次元を一つずつ加えて作れる。
この $L$ への $X,X'$ の制限は比例するため、像が $L$ に入る行列について二つの跡直交式は一つの非零式となる。従って $S^c$ 上で $\operatorname{im}Y=L$ の根 $y$ がある。それを取り、同じ行列像条件の
二空間から、全スカラーが非零な独立な $y'$ を選ぶ。
各スカラー座標汎関数は行列側で可逆な対向座標行列との跡対合である
（$V,W$ の直交と対向座標の非等方性を使う）。よってこの二空間全体で消える座標はない。
$yC$ に対する $y'$ の右コロンは $J_{S^c,k^2}$ となる。
行列の像条件は自動で、スカラーは $S^c$ の外で零となる必要があるからである。
同じ商の次元計算で $(y,y')C=J_{\mathrm{all},L}$ となる。

スカラー台が三点のときは必ず分割の一対を含む。その対から一座標ずつ増やして全体まで進む。
$J_{S,k^2}$ に座標 $j\notin S$ の一次元（スカラーが $j$ にだけ台を持つ元）を加えたとき、
右コロンは $J_{\mathrm{all}\setminus\{j\},k^2}$ である。
次元を各次数で一つずつ増やすので、得るイデアルはちょうど $J_{S\cup\{j\},k^2}$ である。
こうして零、根の主イデアル、二点・三点・全点のスカラー台のイデアル、上記の行列像直線のイデアルが
Koszulフィルトレーションをなす。すなわち各非零イデアルは同じ集合の小さいイデアルに
一次元を一つ加えたもので、そのコロンも集合に属する。
この条件が頂点単純の線型解消を与えることは、各商の短完全列のTor長完全列で従う。
ホモロジー次数に関する帰納の内側で一次生成元数に関して帰納すれば、
コロンの商は一つ低いホモロジー次数しか必要としない。
全正次数イデアルによる商もこの集合にあるから $C$ はKoszulであり、特に二次表示である。
\path{computations/verify_half_period_semisimple.py} は二つの始点の根・主イデアル・
二生成イデアルの各乗法階数を単根・二重根の両方で検算する。
\end{proof}

\begin{proof}[全重みへの適用]
$B^!=C$ が二次関係から従う。両位置の二次空間 $H$ には同じ元
$z=(\omega_1,\ldots,\omega_4,I_2)$ が入る。$D$ の中心単元であり、二周期で $C$ の中心非零因子である。
$B$ の七関係のうち $A_\omega$ に残るものは、$L_\pm$ と、$p_i$ の係数を $\omega$ に対して
直交させた $I$ である。従って直交空間を取れば $A_\omega^!=C/(z)$ である。
中心正則二次式によるこの商もKoszulである。実際 $C\to C/(z)$ のTor環変更スペクトル列は二行しかなく、
下段のホモロジー次数 $p$ のTorから出る可能な微分は上段の $p-2$ のTorの内部次数2ずらしへのものだけである。
その核は極限項で $C$ のKoszul性により内部次数 $p$ に集中する。像も $p$ に関する帰納で同じ次数に集中する。

商のHilbert関数は $1,4,6,4,1$ となり、四次 $D/zH$ の汎関数を $a\mapsto\ell(z^{-1}a)$ とする。
二次 $H/kz$ 上の積対合の根基は零である。$H$ の直交は $k1$ なので割る前の根基が $kz$ だからである。
一次と反対始点の三次 $D/zV_1$ でも $V_0^\perp=V_1$ から完全対合を得る。
従って $A_\omega$ は標準Koszul左右AS正則である。$B$ もKoszulで、
$C$ の級数を符号反転して逆に取れば $B$ は前節の所定の級数を持つ。短完全列の次元比較で $q$ は非零因子である。
具体例 $G=I_4$ でも全ての非零条件が成立し、開条件全体が枠空間に稠密である。
\end{proof}
